% 作者题证的编辑转录副本，不是独立下载的上游原件。
% 来源：https://terrytao.wordpress.com/2009/03/07/infinite-fields-finite-fields-and-the-ax-grothendieck-theorem/

\begin{theorem}[1, Ax--Grothendieck theorem in ${\bf C}$]
Let $P:{\bf C}^n\to{\bf C}^n$ be a polynomial. If $P$ is injective, then $P$ is bijective.
\end{theorem}
\paragraph{1. Proof via finite fields.}
The first observation is that the theorem is utterly trivial in the finite field case:
\begin{theorem}[2, Ax--Grothendieck theorem in $F$]
Let $F$ be a finite field, and let $P:F^n\to F^n$ be a polynomial. If $P$ is injective, then $P$ is bijective.
\end{theorem}
\begin{proof}
Any injection from a finite set to itself is necessarily bijective. (The hypothesis that $P$ is a polynomial is not needed at this stage, but becomes crucial later on.)
\end{proof}
Next, we pass from a finite field $F$ to its algebraic closure $\overline F$.
\begin{theorem}[3, Ax--Grothendieck theorem in $\overline F$]
Let $F$ be a finite field, let $\overline F$ be its algebraic closure, and let $P:\overline F^n\to\overline F^n$ be a polynomial. If $P$ is injective, then $P$ is bijective.
\end{theorem}
\begin{proof}
Our main tool here is Hilbert's nullstellensatz, which we interpret here as an assertion that if an algebraic problem is insoluble, then there exists a finitary algebraic obstruction that witnesses this lack of solution. Specifically, suppose for contradiction that we can find a polynomial $P:\overline F^n\to\overline F^n$ which is injective but not surjective. Injectivity of $P$ means that the algebraic system
\[P(x)=P(y);\qquad x\ne y\]
has no solution over the algebraically closed field $\overline F$; by the nullstellensatz, this implies that for each $j=1,\dots,n$, there must exist an algebraic identity of the form
\[(P(x)-P(y))\cdot Q_j(x,y)=(x_j-y_j)^r \tag{1}\]
for some $r\ge1$ and some polynomial $Q:\overline F^n\times\overline F^n\to\overline F^n$ that specifically witnesses this lack of solvability. Similarly, lack of surjectivity means the existence of an $z_0\in\overline F^n$ such that the algebraic system
\[P(x)=z_0\]
has no solution over the algebraically closed field $\overline F$; by another application of the nullstellensatz, there must exist an algebraic identity of the form
\[(P(x)-z_0)\cdot R(x)=1 \tag{2}\]
for some polynomial $R:\overline F^n\to\overline F^n$ that specifically witnesses this lack of solvability.

Fix $Q_j,z_0,R$ as above, and let $k$ be the subfield of $\overline F$ generated by $F$ and the coefficients of $P,Q_j,z_0,R$. Then we observe (thanks to our explicit witnesses (1), (2)) that the counterexample $P$ descends from $\overline F$ to $k$; $P$ is a polynomial from $k^n$ to $k^n$ which is injective but not surjective. But $k$ is finitely generated, and every element of $k$ is algebraic over the finite field $F$, thus $k$ is finite. But this contradicts Theorem 2.
\end{proof}
\begin{proof}[The complex case; the author's continuation]
The complex case ${\bf C}$ follows by a slight extension of the argument used to prove Theorem 3. Indeed, suppose for contradiction that there is a polynomial $P:{\bf C}^n\to{\bf C}^n$ which is injective but not surjective. As ${\bf C}$ is algebraically closed, we may invoke the nullstellensatz as before and find witnesses (1), (2) for some $Q,z_0,R$.

Now let $k={\bf Q}[\mathcal C]$ be the subfield of ${\bf C}$ generated by the rationals ${\bf Q}$ and the coefficients $\mathcal C$ of $P,Q,z_0,R$. Then we can descend the counterexample to $k$. This time, $k$ is not finite, but we can descend it to a finite field (and obtain the desired contradiction) by a number of methods.

One approach, which is the one taken by Serre, is to quotient the ring ${\bf Z}[\mathcal C]$ generated by the above coefficients by a maximal ideal, observing that this quotient is necessarily a finite field.
\end{proof}
